\section{From a guitar in a sandstorm to moment asymptotics}\label{S:reader-history}

Imagine a guitar left outside during a sandstorm. Grains strike a string at
unpredictable places and times, and the string carries each disturbance away.
Walsh used this picture in his 1984 Saint-Flour lectures, published in
1986~\cite{walsh:86:introduction}*{p.~282, (1.14)}. His worked example idealizes
the string as infinite, initially at rest, and driven by additive white noise.
The integral formulation makes sense of motion too irregular for the classical
derivatives in the wave equation. Forty years since that landmark exposition,
the questions extend from constructing the motion to measuring its growth.

\begin{figure}
  \centering
  \begin{tikzpicture}
    \begin{scope}
      \clip[rounded corners=9pt]
      (-0.24\textwidth,-0.24\textwidth) rectangle
      (0.24\textwidth,0.24\textwidth);
      \node[inner sep=0pt,outer sep=0pt] (art) {
      \includegraphics[width=0.48\textwidth]{media/guitar-sandstorm/illustration.png}};
      % Soft white rings blend the image into the paper without changing the art.
      \begin{scope}
        \clip (art.south west) rectangle (art.north east);
        \foreach \k in {24,23,...,1} {
        \draw[white,opacity=0.12,line width=\k*0.65pt,
        rounded corners=9pt]
        (art.south west) rectangle (art.north east);
        }
      \end{scope}
      % The outer white frame also removes the square corner remnants.
      \draw[white,line width=3pt,rounded corners=9pt]
      (art.south west) rectangle (art.north east);
    \end{scope}
  \end{tikzpicture}
  \caption[A guitar in a sandstorm]{AI-generated illustration inspired by
    Walsh's guitar-in-a-sandstorm example. Artwork from Le Chen's
    \href{https://github.com/SPDEngine/SPDE-Visuals}{SPDE-Visuals collection},
    shared under the
    \href{https://creativecommons.org/publicdomain/zero/1.0/}{CC0 public-domain dedication}.
    This is an artistic illustration, not a
    numerical simulation. Walsh's worked example uses an infinite string to
    avoid boundary conditions. Companion videos and source information are
    included in \path{media/guitar-sandstorm/}.}
  \label{Fig:walsh-guitar}
\end{figure}


Dalang's integration theory~\cite{dalang:99:extending} widened the treatment of
singular Green kernels and spatial covariance. Probabilistic moment
methods~\cite{dalang.mueller.ea:08:feynman-kac-type} and explicit
one-dimensional second-moment formulas~\cite{chen.dalang:15:moment} brought the
long-time fluctuations into focus. Exact asymptotics for static Skorohod
noise~\cite{balan.chen.ea:22:exact} and static Stratonovich
noise~\cite{chen.hu:25:hyperbolic} then sharpened the question: what is the
leading constant, and how does the meaning of the noise product affect it?

The present paper asks this question for the second moment with fractional
temporal correlation and Stratonovich multiplication. We give an upper bound at
the correct scale and a variational expression for the second-moment limit. The
working-draft notice records the review status. This account follows selected
developments since Walsh's lectures; it does not cover the earlier history or
all later work. The book's account in \TNref{Section}{S:wave-history} and its
literature comparison retain the hypotheses and distinctions behind these
milestones.
